% !TeX root = ../colloquium-2.tex

% Определение дифференцируемой функции. 
% Определение дифференциала.
% Линейная функция.

\phantomsection\label{sec_Definition05}
\begin{definition}[Определение дифференцируемой функции. 
                   Определение дифференциала. 
                   Линейная функция.]
    \subsection*{Определение дифференцируемой функции.}
    \begin{note}[purple]
        Функция $f$, определённая в некоторой окрестности точки $a$,
        называется \textbf{дифференцируемой} в точке $a$, если 
        существуют такие число $A$ и функция $\alpha$, что 
        для всех $h$ из некоторой проколотой окрестности нуля 
        выполнено равенство 
        $$f(a + h) - f(a) = Ah + \alpha(h)h,$$ 
        где $\displaystyle{\lim_{h \rightarrow 0} \alpha(h) = 0}$. \\
        При этом $A$ и $\alpha$ зависят и от точки $a$, поэтому часто
        равенство записывают в виде 
        $$ f(a + h) - f(a) = A(a)h + \alpha(a, h)h $$
    \end{note}

    \subsection*{Определение дифференциала.}
    \begin{note}[purple]
        Функция $h \mapsto Ah$ называется \textbf{дифференциалом} 
        функции $f$ в точке $a$. Она обозначается $df(a)$ или
        $df | _{x = a}$, то есть $df(a)(h) = df(h)|_{x = a} = Ah$.
    \end{note}

    \subsection*{Линейная функция.}
    \begin{note}[purple]
        Линейной функцией называется функция вида $f = kx + b$, 
        где $x$ - независимая переменная, $k$ и $b$ - числа. 
    \end{note}
\end{definition}